%% =====================================================================
%% 浙江工商大学本科毕业论文（设计）Quarto/LaTeX 样式文件
%% 依据：《浙江工商大学正文格式要求20260427.docx》
%% 本文件只管格式，不写内容。学生一般不需要改这里。
%% =====================================================================

\usepackage{geometry}
\usepackage{graphicx}
\usepackage{etoolbox}
\usepackage{zhnumber}
\usepackage{titlesec}
\usepackage{tocloft}
\usepackage{caption}
\usepackage{fancyhdr}
\usepackage{pifont}
\usepackage[perpage]{footmisc}
\usepackage{longtable}
\usepackage{booktabs}
\usepackage{framed}
\usepackage{enumitem}
\setlist[itemize]{itemsep=0pt,topsep=0pt,parsep=0pt,leftmargin=2em}

%% ---------- 1. 中文字体：按操作系统自动降级 ----------
%% 2026-11-07：复位 Quarto 模板默认的 Scale=MatchLowercase，字号严格等于 \fontsize
\defaultfontfeatures{Scale=1}
%% Windows: SimSun / SimHei；macOS: Songti SC / Heiti SC；
%% Linux（本机）: Noto Serif/Sans CJK SC。找不到就退回 Fandol。
\IfFontExistsTF{SimSun}
  {\setCJKmainfont{SimSun}[AutoFakeBold=2.5]}
  {\IfFontExistsTF{Songti SC}
    {\setCJKmainfont{Songti SC}[AutoFakeBold=2.5]}
    {\IfFontExistsTF{Noto Serif CJK SC}
      {\setCJKmainfont{Noto Serif CJK SC}[AutoFakeBold=2.5]}
      {\setCJKmainfont{FandolSong}}}}

\IfFontExistsTF{SimHei}
  {\setCJKsansfont{SimHei}}
  {\IfFontExistsTF{Heiti SC}
    {\setCJKsansfont{Heiti SC}}
    {\IfFontExistsTF{Noto Sans CJK SC}
      {\setCJKsansfont{Noto Sans CJK SC}}
      {\setCJKsansfont{FandolHei}}}}

%% 西文与阿拉伯数字：Times New Roman（找不到用 Times 等价字体）
\IfFontExistsTF{Times New Roman}
  {\setmainfont{Times New Roman}[Scale=1]}
  {\setmainfont{TeX Gyre Termes}[Scale=1]}

%% 日期里的“○”按中文全角圆排（2026-11-10）
%% 封面日期写作“二○二六年十一月十一日”，其中的圈是 WHITE CIRCLE U+25CB，不在
%% xeCJK 默认中文区段内，会被 Times 接走：字形窄一圈，左右各多出约 4 磅中西文间距，
%% 印出来像“二 ○ 二”。归入 CJK 类后交宋体排，宽度与“〇”一致（16pt 下 15.94 磅）。
\xeCJKDeclareCharClass{CJK}{"25CB}

%% 正文用宋体、标题用黑体
\providecommand{\songti}{}
\providecommand{\heiti}{}
\renewcommand{\songti}{\rmfamily}
\renewcommand{\heiti}{\sffamily}

%% ---------- 2. 字体号别名（中文习惯） ----------
\newcommand{\xiaosi}{\fontsize{12pt}{20pt}\selectfont}   % 小四，固定行距 20 磅
\newcommand{\sihao}{\fontsize{14pt}{28pt}\selectfont}    % 四号，2 倍行距
\newcommand{\sanhao}{\fontsize{16pt}{20pt}\selectfont}   % 三号
\newcommand{\wuhao}{\fontsize{10.5pt}{16pt}\selectfont}  % 五号
\newcommand{\xiaowu}{\fontsize{9pt}{11pt}\selectfont}    % 小五

%% ---------- 3. 正文行距：固定值 20 磅 ----------
\renewcommand{\normalsize}{\fontsize{12pt}{20pt}\selectfont}
\setlength{\parindent}{2em}         % 首行缩进两字符
\linespread{1}
%% 2026-11-10：与大论文模板一致，段间距归零；否则封面每段会多出约 9.4 磅
\setlength{\parskip}{0pt}

%% ---------- 4. 标题编号与格式 ----------
%% 一级 一、 四号宋体加粗，居中，上下各空一行，2 倍行距
%% 二级 （一） 小四宋体，左对齐空两格，不空行
%% 三级 1. 小四宋体
%% 四级 （1） 小四宋体，与正文连排
\renewcommand{\thesection}{\zhnum{section}、}
\renewcommand{\thesubsection}{（\zhnum{subsection}）}
\renewcommand{\thesubsubsection}{\arabic{subsubsection}.}
\renewcommand{\theparagraph}{（\arabic{paragraph}）}

\titleformat{\section}[block]
  {\centering\songti\fontsize{14pt}{28pt}\selectfont\bfseries}
  {\thesection}{0.3em}{}
\titlespacing{\section}{0pt}{20pt}{20pt}

\titleformat{\subsection}[block]
  {\songti\fontsize{12pt}{20pt}\selectfont}
  {\thesubsection}{0.3em}{}
\titlespacing{\subsection}{2em}{0pt}{0pt}

\titleformat{\subsubsection}[block]
  {\songti\fontsize{12pt}{20pt}\selectfont}
  {\thesubsubsection}{0.3em}{}
\titlespacing{\subsubsection}{2em}{0pt}{0pt}

\titleformat{\paragraph}[runin]
  {\songti\fontsize{12pt}{20pt}\selectfont}
  {\theparagraph}{0em}{}
\titlespacing*{\paragraph}{2em}{0pt}{0.5em}

%% ---------- 5. 目录：到二级，小四宋体，固定 20 磅 ----------
\AtBeginDocument{%
  \renewcommand{\contentsname}{目\hspace{2em}录}%
  \setcounter{tocdepth}{2}%
}
\renewcommand{\cfttoctitlefont}{\hspace*{\fill}\heiti\fontsize{16pt}{20pt}\selectfont}
\renewcommand{\cftaftertoctitle}{\hspace*{\fill}\par}
\renewcommand{\cftsecfont}{\songti\fontsize{12pt}{20pt}\selectfont}
\renewcommand{\cftsecpagefont}{\fontsize{12pt}{20pt}\selectfont}
\renewcommand{\cftsubsecfont}{\songti\fontsize{12pt}{20pt}\selectfont}
\renewcommand{\cftsubsecpagefont}{\fontsize{12pt}{20pt}\selectfont}
\setlength{\cftsecnumwidth}{36pt}
\setlength{\cftsubsecnumwidth}{42pt}
\setlength{\cftsubsecindent}{30.7pt}

%% ---------- 6. 图表题注：五号宋体加粗，图下图题、表上表题 ----------
\DeclareCaptionFont{wuhao}{\fontsize{10.5pt}{16pt}\selectfont}
\captionsetup{
  font=wuhao,
  labelfont={bf},
  textfont={bf},
  labelsep=space,
  labelformat=simple,
  singlelinecheck=true
}
%% 中文图表名放到 \AtBeginDocument，避免被 babel/Quarto 模板覆盖
%% 题注标签格式改为“图1”“表1”（校规不加空格），标签与题名之间空一格
\DeclareCaptionLabelFormat{zhnum}{#1#2}
\captionsetup{labelformat=zhnum}
\AtBeginDocument{\renewcommand{\figurename}{图}\renewcommand{\tablename}{表}}

%% ---------- 7. 三线表：上下 1.5 磅，中间 0.75 磅 ----------
\setlength{\heavyrulewidth}{1.5pt}
\setlength{\lightrulewidth}{0.75pt}
\setlength{\cmidrulewidth}{0.75pt}
\setlength{\aboverulesep}{0pt}
\setlength{\belowrulesep}{0pt}
\AtBeginEnvironment{longtable}{\fontsize{10.5pt}{16pt}\selectfont}
\AtBeginEnvironment{tabular}{\fontsize{10.5pt}{16pt}\selectfont}

%% ---------- 8. 页边距（与正文格式要求一致） ----------
\geometry{top=2.54cm,bottom=2.54cm,left=3.17cm,right=3.17cm,headheight=13pt,headsep=0.45cm,footskip=0.655cm}

%% 公式：居中排版，编号右对齐（与论文正文模板一致，classoption 不再用 fleqn）

%% ---------- 9. 页眉页脚 ----------
%% 页眉：两行小五(9pt)宋体、左对齐；横线按 Word 默认 0.75 磅（校规与范例均无明文，
%% 2026-11-07 主人指令）。
%% 第一行标明文档类型，第二行“毕业论文（设计）题目：”后接题目。
%% 2026-11-10：封面标题按教务处 doc 原件改为“本科毕业论文开题报告”（去掉“（设计）”），
%% 页眉第一行同步改为“毕业论文：开题报告”，免得封面与页眉两个说法。
%% 封面用 \thispagestyle{empty} 不带页眉页脚。
\providecommand{\zjgsuTitle}{}
\pagestyle{fancy}
\makeatletter\let\ps@plain\ps@fancy\makeatother
\fancyhf{}
\renewcommand{\headrulewidth}{0.75pt}
\fancyhead[L]{\songti\fontsize{9pt}{13pt}\selectfont
  毕业论文：开题报告\\毕业论文（设计）题目：\zjgsuTitle}
\fancyfoot[C]{\fontsize{9pt}{9pt}\selectfont -\ \thepage\ -}

%% ---------- 10. 脚注：五号宋体，每页重新编号 ①②③ ----------
\renewcommand{\thefootnote}{\ding{\numexpr171+\value{footnote}}}
\renewcommand{\footnotesize}{\fontsize{10.5pt}{16pt}\selectfont}
\AtBeginEnvironment{footnotesize}{\xiaowu}

%% ---------- 11. 参考文献：五号，悬挂缩进由 CSL 的 hanging-indent 控制 ----------
\renewcommand{\refname}{参考文献}

%% ---------- 12. 超链接不显示颜色 ----------
\AtBeginDocument{%
  \hypersetup{
    colorlinks=true,
    linkcolor=black,
    citecolor=black,
    urlcolor=black,
    filecolor=black
  }%
}

%% ---------- 13. 开题报告的栏框（可跨页） ----------
%% 2026-11-07：字号复位 Scale 后首框题目行仅差约 1 磅溢出，栏框内边距 3 → 2 磅即收进一行
\setlength{\FrameSep}{2.5pt}
\newenvironment{zjgsuBox}{\begin{framed}}{\end{framed}}
